% Generated draft: requires mathematical and source review.
\documentclass[11pt]{article}
\usepackage[margin=25mm]{geometry}
\usepackage{fontspec}
\setmainfont{Latin Modern Roman}
\usepackage{amsmath,amssymb,mathtools,unicode-math}
\setmathfont{Latin Modern Math}
\usepackage{enumitem,needspace,xurl,hyperref}
\setlength{\parindent}{0pt}
\setlength{\parskip}{5pt}
\setlength{\emergencystretch}{3em}
\newcommand{\N}{\mathbb{N}}
\newcommand{\Z}{\mathbb{Z}}
\newcommand{\Q}{\mathbb{Q}}
\newcommand{\R}{\mathbb{R}}
\newcommand{\C}{\mathbb{C}}
\newcommand{\F}{\mathbb{F}}
\newcommand{\A}{\mathbb{A}}
\newcommand{\PP}{\mathbb{P}}
\newcommand{\E}{\mathbb{E}}
\newcommand{\eps}{\varepsilon}
\newcommand{\dd}{\,\mathrm{d}}
\begin{document}
\section*{UnsolvedMath ID: 20001893}
\textbf{Source problem code:} catalogue entry
\subsection{Definitions and mathematical statement}
Let \((M^n,g)\) be closed with negative sectional curvature, and let \(\phi^t\) be its geodesic flow. At an Oseledets-regular point write the positive exponents as
\(0<\lambda_{\min}\leq\cdots\leq\lambda_{\max}\).
Question 5.4.2 asks whether some metric has
\[
 \lambda_{\max}/\lambda_{\min}>2
\]
for typical geodesics, rather than merely exceptional periodic ones. The paired aid interprets typical as almost every point for normalized Liouville measure. This convention is not explicit in the original wording. A surface has only one positive exponent, so this ratio question is nontrivial only in higher dimension.

The aid's obstruction assumes
\[
 -b^2\leq K\leq-a^2<0,\qquad0<a\leq b.
\]
It reports that every positive Lyapunov exponent, for every invariant probability measure and every Oseledets-regular point, lies in \([a,b]\). Hence
\[
 \lambda_{\max}/\lambda_{\min}\leq b/a,
\]
and a typical ratio greater than \(2\) requires \(b/a>2\). It is impossible under weak quarter pinching \(b/a\leq2\), including equality. This is a necessary condition, not an example outside the pinching range.

For its occupation refinement, let \(J(t)\) be an unstable covariant Jacobi direction with exponent \(\lambda\), let \(e(t)=J(t)/\|J(t)\|\), and let \(U^u(t)\) be the unstable Riccati operator. The instantaneous rate is
\[
 \rho(t)=\langle U^u(t)e(t),e(t)\rangle\in[a,b],\qquad
 \lambda=\lim_{T\to\infty}\frac1T\int_0^T\rho(t)\,dt.
\]
For \(a<c<b\), the aid reports
\[
 \lambda>c\ \Longrightarrow\
 \liminf_{T\to\infty}\frac{|\{t\in[0,T]:\rho(t)\geq c\}|}{T}
 \geq\frac{\lambda-c}{b-c},
\]
\[
 \lambda<c\ \Longrightarrow\
 \liminf_{T\to\infty}\frac{|\{t\in[0,T]:\rho(t)\leq c\}|}{T}
 \geq\frac{c-\lambda}{c-a}.
\]
Time-set size is Lebesgue measure; the bounds concern the chosen covariant directions, not just extremal sectional curvatures. The report says its literature check through 2026 found no typical-ratio-above-two example and leaves the existence question open.

The concatenated source includes additional regularity context. It attributes smoothness of the Anosov splitting precisely for locally symmetric metrics to [Be-Fo-La1992] and [Be-Co-Ga1995]. For surfaces it attributes to Hurder--Katok [Hur-Kat1990] a smoothness implication from a stable or unstable foliation of regularity written as \(C^{1+o(x\log|x|)}\). The precise modulus convention is not supplied. The source describes the higher-dimensional implication \(C^2\Rightarrow C^\infty\) as expected and then open; no new status check of that separate problem is provided.

The tail also begins Section 5.5, ``How many closed geodesics of length \(\leq T\) exist?'', communicated by Pollicott. It attributes to [Margulis2004], for closed negatively curved \(M\), the asymptotic
\[
 N(T)\sim\frac{e^{hT}}{hT}.
\]
Here \(h\) is understood as the geodesic flow's topological entropy; primitive-orbit and orientation counting conventions are omitted. This adjacent counting fragment and its reference are retained separately from the Lyapunov existence question, without inventing a missing refinement.
\subsection{Short English statement}
The supplied aid bounds every positive exponent by curvature pinching, ruling out a ratio above two under weak quarter pinching. It also gives quantitative necessary occupation bounds for unstable Riccati rates and still reports the typical-ratio existence question open. The source's concatenated splitting-regularity and closed-geodesic-counting material is retained with its own qualifications.
\subsection{Sources}
\begin{itemize}
\item Local source record: catalogue entry, catalogue entry; canonical \texttt{aim-geometry-notes.json}, \texttt{notes[230]}, source item 5.4, extracted Question 5.4.2. Title: ``A pinching and Riccati-occupation obstruction to typical Lyapunov spread beyond two.''
\item Keith Burns and Vladimir Matveev, \emph{Open problems and questions about geodesics}, AIM workshop ``Geodesics,'' supplied reference \url{https://aimath.org/pastworkshops/geodesicsproblems.pdf}. This draft uses the supplied record without accessing a PDF.
\item Supplied source references [Be-Fo-La1992], [Be-Co-Ga1995], Hurder--Katok [Hur-Kat1990], and [Margulis2004]. The counting section is communicated by Pollicott. No full bibliographic entries are supplied.
\item Research attempt 1 supplies the pinching and Riccati-occupation obstruction and an existence-status check through 2026. Registry dated 14 August 2026: ``partially\_solved,'' ``PARTIAL-PROGRESS,'' statement status ``exact''; attempt partial, research-run status open, rigor ``proved,'' novelty confidence low. The machine-generated, non-peer-reviewed aid requires independent verification before citation.
\end{itemize}
\end{document}
